% Generated by scripts/assemble.py; corrections belong in source/corrections.json.
\examhead{107}{合肥工业大学 2024 年数学分析试卷}{105}
% p106-b21; source PDF page 106
\par\Needspace{6\baselineskip}
\blockanchor{p106-b21}
\question{1}(10 分) 求极限\par
\nopagebreak[4]
% p106-b22; source PDF page 106
\blockanchor{p106-b22}
\displayformula{\lim _ {n \rightarrow \infty} n ^ {2} \left(3 ^ {\frac {1}{n}} - 3 ^ {\frac {1}{n + 1}}\right).}
% p106-b23; source PDF page 106
\par\Needspace{6\baselineskip}
\blockanchor{p106-b23}
\question{2}(10 分) 求极限\par
\nopagebreak[4]
% p106-b24; source PDF page 106
\blockanchor{p106-b24}
\displayformula{\lim _ {x \to 0} \frac {1}{x ^ {4}} \left[ \ln (1 + \sin^ {2} x) - 6 (\sqrt [ 3 ]{2 - \cos x} - 1) \right].}
% p106-b25; source PDF page 106
\blockanchor{p106-b25}
\question{3}(10 分) 判断数项级数 $\sum_{n=2}^{\infty} (-1)^{n} \ln \left(1 + \frac{1}{n}\right)$ 是否收敛, 如果收敛, 则判断是条件收敛还是绝对收敛.\par
% p107-b00; source PDF page 107
\par\Needspace{6\baselineskip}
\blockanchor{p107-b00}
\question{4}(15 分) 求极限\par
\nopagebreak[4]
% p107-b01; source PDF page 107
\blockanchor{p107-b01}
\displayformula{\lim _ {n \rightarrow \infty} \frac {\ln n}{\ln \left(1 ^ {2023} + 2 ^ {2023} + \cdots + n ^ {2023}\right)}.}
% p107-b02; source PDF page 107
\par\Needspace{6\baselineskip}
\blockanchor{p107-b02}
\question{5}(15 分) 设 $f(x)$ 在 $[a, b]$ 上可积, 且在 $x = b$ 处连续, 证明:\par
\nopagebreak[4]
% p107-b03; source PDF page 107
\blockanchor{p107-b03}
\displayformula{\lim _ {n \rightarrow \infty} \frac {n + 1}{(b - a) ^ {n + 1}} \int_ {a} ^ {b} (x - a) ^ {n} f (x) \mathrm{d} x = f (b).}
% p107-b04; source PDF page 107
\blockanchor{p107-b04}
\question{6}(15 分) 设函数 $f(x)$ 在 $(- \infty, +\infty)$ 上有界且 $f''(x) \geq 0$. 证明: $f(x)$ 为常值函数.\par
% p107-b05; source PDF page 107
\blockanchor{p107-b05}
\question{7}(15 分) 求幂级数 $\sum_{n=0}^{\infty}(n+2)(n+1)x^{n}$ 的收敛域及和函数.\par
% p107-b06; source PDF page 107
\blockanchor{p107-b06}
\question{8}(15 分) 证明: $F(x) = \int_{1}^{+\infty} \frac{y \sin xy}{1 + y^2} \, \mathrm{d}y$ 在 $(0, +\infty)$ 上连续.\par
% p107-b07; source PDF page 107
\blockanchor{p107-b07}
\question{9}(15 分) 设有方程 $\cos(x + y) + e^{y + x^2} - x^3 y^3 = 2$.\par
% p107-b08; source PDF page 107
\blockanchor{p107-b08}
(1) 证明: 在点 $(0,0)$ 的某邻域内能确定唯一的隐函数 $y = f(x)$.\par
% p107-b09; source PDF page 107
\blockanchor{p107-b09}
(2) 判断在  $(0,0)$  的某邻域内能否确定隐函数  $x = f(y)$ ? 并说明理由.\par
% p107-b10; source PDF page 107
\par\Needspace{6\baselineskip}
\blockanchor{p107-b10}
\question{10}(15 分) 已知\par
\nopagebreak[4]
% p107-b11; source PDF page 107
\blockanchor{p107-b11}
\displayformula{f (x, y) = \left\{ \begin{array}{l l} x y \frac {x ^ {2} - y ^ {2}}{x ^ {2} + y ^ {2}}, & (x, y) \neq (0, 0); \\ 0, & (x, y) = (0, 0). \end{array} \right.}
% p107-b12; source PDF page 107
\blockanchor{p107-b12}
证明： $f(x,y)$ 在 $(0,0)$ 处连续，并计算 $f_{x}(0,0)$ 和 $f_{xy}(0,0)$\par
% p107-b13; source PDF page 107
\par\Needspace{6\baselineskip}
\blockanchor{p107-b13}
\question{11}(15 分) 求曲面积分\par
\nopagebreak[4]
% p107-b14; source PDF page 107
\blockanchor{p107-b14}
\displayformula{\iint_ {S} x \mathrm{d} y \mathrm{d} z + y \mathrm{d} z \mathrm{d} x + [ (z + 1) ^ {2} + e ^ {- (x ^ {2} + y ^ {2})} ] \mathrm{d} x \mathrm{d} y.}
% p107-b15; source PDF page 107
\blockanchor{p107-b15}
其中 S 为  $z = 1 - x^{2} - y^{2} (0 \leq z \leq 1)$ ，方向取下侧.\par
